\section{Metodologia}

Este capítulo detalha os procedimentos computacionais e o planejamento estatístico adotados para a construção, otimização e avaliação automatizada do sistema RAG aplicado ao acervo normativo do MPES. A pesquisa caracteriza-se como um estudo experimental empírico e quantitativo, estruturado rigorosamente para mensurar o impacto isolado e as interações sinérgicas de diferentes hiperparâmetros arquiteturais na qualidade da informação recuperada e gerada.

Para garantir a reprodutibilidade e a validade interna, a exposição metodológica segue uma ordem lógico-construtiva: inicia-se com a formalização do Delineamento Experimental e a definição das variáveis. Em seguida, detalha-se a engenharia de dados e o processo de curadoria sintética do conjunto de teste. Por fim, descrevem-se os protocolos estritos de execução em lote das inferências e a modelagem estatística não-paramétrica ART empregada para testar as hipóteses do estudo.
